%---------- Inleiding ---------------------------------------------------------

% TODO: Is dit voorstel gebaseerd op een paper van Research Methods die je
% vorig jaar hebt ingediend? Heb je daarbij eventueel samengewerkt met een
% andere student?
% Zo ja, haal dan de tekst hieronder uit commentaar en pas aan.

%\paragraph{Opmerking}

% Dit voorstel is gebaseerd op het onderzoeksvoorstel dat werd geschreven in het
% kader van het vak Research Methods dat ik (vorig/dit) academiejaar heb
% uitgewerkt (met medesturent VOORNAAM NAAM als mede-auteur).
% 

\section{Inleiding}%
\label{sec:inleiding}

Steden in het Noordzeegebied worden in toenemende mate geconfronteerd met risico's op zowel wateroverlast als temperatuurstijging 
als gevolg van klimaatverandering \autocite{EEA2020}.
Voor de stad Gent vormt hitteoverlast een aantoonbaar klimaatgerelateerd risico.
Binnen Europese vergelijkingen behoort Gent tot de steden met een relatief hoge kwetsbaarheid voor hitte \autocite{Tapia2017}.

Stedelijke gebieden zijn bijzonder gevoelig voor temperatuurstijgingen door het stedelijk hitte-eiland effect.
Deze Urban heat islands (UHI) worden gedefinieerd als het fenomeen waarbij stedelijke gebieden 
hogere temperaturen vertonen dan hun omliggende landelijke gebieden \autocite{Phelan2015}.
Volgens \textcite{Zheng2026} kan het temperatuurverschil oplopen tot 12°C in extreme gevallen.
Deze studie identificeert stedelijke vergroening als één van de factoren die samenhangen met lagere oppervlaktetemperaturen.

De stedelijke vergroening in Gent omvat onder meer grootschalige boomaanplanting.
Hoewel verschillende studies aantonen dat vergroening een positief effect heeft op stedelijk klimaat \autocite{Knight2021, Zheng2026, Phelan2015}, 
is het kwantitatieve verband tussen boomdensiteit en temperatuur afhankelijk van lokale ruimtelijke kenmerken \autocite{Wang2024}.
Voor deze reden is het effect van boomplanting moeilijk te beoordelen zonder specifieke onderzoek.
In Gent ontbreekt een recente, gedetailleerde ruimtelijke analyse die het effect van boomdensiteit op temperatuur 
kwantificeert rekening houdend met lokale stedelijke kenmerken.

De voorgestelde bachelorproef heeft als doel deze relatie tussen vergroening en temperatuur in Gent te onderzoeken 
aan de hand van data-analyse.
Daarnaast kunnen op basis van de analyse prioriteitsgebieden worden geïdentificeerd waar vergroening het meest 
effectief zou zijn voor temperatuurreductie.
De resultaten kunnen bijdragen aan een beter onderbouwde evaluatie van stedelijke vergroening als klimaatadaptatiemaatregel.

%---------- Stand van zaken ---------------------------------------------------

\section{Literatuurstudie}%
\label{sec:literatuurstudie}


% Voor literatuurverwijzingen zijn er twee belangrijke commando's:
% \autocite{KEY} => (Auteur, jaartal) Gebruik dit als de naam van de auteur
%   geen onderdeel is van de zin.
% \textcite{KEY} => Auteur (jaartal)  Gebruik dit als de auteursnaam wel een
%   functie heeft in de zin (bv. ``Uit onderzoek door Doll & Hill (1954) bleek
%   ...'')

\subsection{Urban Heat Islands en temperatuurmetingen}

Het stedelijke hitte-eilandeffect (UHI) beschrijft het fenomeen waarbij 
stedelijke gebieden hogere temperaturen vertonen dan omliggende 
landelijke gebieden, met temperatuurverschillen tot 12°C in extreme 
gevallen \autocite{Phelan2015, Zheng2026}. Dit effect wordt veroorzaakt 
door verharding, verminderde verdamping en afwezigheid van vegetatie.

Om UHI-effecten te kwantificeren worden verschillende meetmethoden 
gebruikt. \textcite{Moncada-Morales2025} identificeerden in een 
bibliometrische analyse van 130 studies drie hoofdmethoden: Land Surface 
Temperature (LST) via satellietbeelden (86 studies), grondgebonden 
luchttemperatuurmetingen (25 studies), of beide gecombineerd (21 studies).

Dit onderzoek maakt gebruik van LST-data vanwege de hoge ruimtelijke 
resolutie (30-100m), die essentieel is om gelokaliseerde effecten van 
individuele bomen en specifieke landgebruikstypen te detecteren. In 
tegenstelling tot grondgebonden weerstations bieden satellietbeelden 
volledige ruimtelijke dekking van het studiegebied.

\subsection{Invloed van stedelijke vergroening op temperatuur}

In \textcite{Knight2021} zijn 308 studies beschouwd om een algemene conclusie te trekken 
rond het effect van stedelijke vergroening op ozon, UV en temperatuur.
Hierin is gevonden dat luchttemperatuur onder bomen gemiddeld 0.8°C koeler was, 
terwijl stedelijke bossen gemiddeld 1.6°C koeler waren.
Deze studie duidt aan dat vergroening een significante impact heeft op temperatuur en menselijke comfort.

Hoewel het koelende effect van vergroening duidelijk is, 
varieert de sterkte ervan sterk afhankelijk van lokale omstandigheden \autocite{Gao2025, Kim2021}. 
Studies gebruiken verschillende onderzoeksmethoden en analysemodellen, 
en dezelfde interventie kan tot verschillende effecten leiden in verschillende stedelijke contexten. 
Dit maakt het moeilijk om algemene voorspellingen te doen die een vergroeningsbeleid kunnen onderbouwen zonder lokale analyse.

Deze context-afhankelijkheid onderstreept de noodzaak van lokaal-specifiek onderzoek, zoals deze voorgestelde studie voor Gent.

%---------- Methodologie ------------------------------------------------------
\section{Methodologie}%
\label{sec:methodologie}

Om het effect van vergroening in Gent te beschouwen wordt de volgende methodologie gebruikt.

\subsection{Bronnen van data}

Voor deze studie zijn 3 soorten data nodig:

\begin{itemize}
    \item Vergroeningsdata: Open Data Portaal Gent (ODP) \autocite{StadGentBomen}
    \item Temperatuurdata:
    \begin{itemize}
        \item LST-data: Online Global Land Surface Temperature Estimation from Landsat \autocite{Parastatidis2017}
        \item Andere: UGent MOCCA
    \end{itemize}
    \item Landgebruiksdata:
    \begin{itemize}
        \item Wegeninformatiesysteem Gent: ODP
        \item Urban Atlas Land Cover/Land Use 2018
    \end{itemize}
\end{itemize}

Gegevens over landgebruik kunnen de onvoorspelbaarheid van lokale weersomstandigheden mitigeren, 
door een beeld te geven van de geometrie en verharding van de omgeving.

\subsection{Fuseren van temperatuurdata}

Zodra temperatuurdata is opgehaald kan de volgende fase gebeuren, namelijk het fuseren van deze data.

LST-data uit \textcite{Parastatidis2017} heeft hoge spatiale resolutie: 30m en 100m, 
maar lage temporale resolutie: 3 huidig actieve satellieten, met datapunten in Gent ongeveer om de week.
Dit gebrek aan temporale resolutie is mogelijk te mitigeren door deze data te fuseren met een andere dataset 
die wel een hoge temporale resolutie heeft, bijvoorbeeld vaste weerstations.
Deze fusie gebeurt aan de hand van ESTARFM, gedefinieerd in \textcite{Zhu2010}.
ESTARFM is ideaal voor deze studie omdat het specifiek ontwikkeld is voor het fuseren van metingen afkomstig van remote sensing.

Na het fuseren van LST data met data van weerstations bevat de resulterende dataset zowel hoge spatiale- en temporale resolutie.
Dat betekent dat er veel meer datapunten zijn om mee te werken, dit zal de nauwkeurigheid van de eventuele analyse verhogen.

\subsection{Data combineren}

Om analyse en voorspelling te doen, worden alle databronnen gecombineerd tot één geïntegreerde dataset.

\subsubsection{Temperatuur- en vergroeningsdata}

De dataset wordt geïnitialiseerd op basis van de gefuseerde LST-data en behoudt dezelfde ruimtelijke resolutie. 
Voor elke cel wordt het aantal bomen berekend, zowel binnen de cel zelf als in omliggende cellen.
Er wordt rekening gehouden met de plantdatum van elke boom, 
zodat het effect van nieuw aangeplante vergroening in de tijd kan worden gevolgd. 
Tevens wordt bijgehouden of een temperatuurmeting afkomstig is van directe LST-observatie of van ESTARFM-fusie.

\subsubsection{Landgebruiksdata}

Per LST-cel worden landgebruikspercentages berekend uit twee databronnen. 
Het Wegeninformatiesysteem Gent levert gedetailleerde straatsamenstelling,
terwijl Urban Atlas 2018 de overige landgebruiksklassen beschrijft.

Beide datasets bevatten een groot aantal landgebruiksklassen. Deze worden 
gehergroepeerd in categorieën met gelijkaardig effect op klimaat 
(zie appendix \ref{app:landgebruik} voor de volledige classificatie):

Gent WIS bevat 2 velden per record:
\begin{itemize}
    \item Bestemming: Gemotoriseerd verkeer, Langzaam verkeer, Groen, Andere
    \item Materiaal: Volledig verhard, Semi-permeabel, Onverhard, Andere
\end{itemize}

Urban Atlas:
\begin{itemize}
    \item Hoog verhard stedelijk
    \item Laag verhard stedelijk 
    \item Transportinfrastructuur
    \item Stedelijk groen
    \item Natuurlijk groen
    \item Water en wetland
    \item Kaal/in ontwikkeling
\end{itemize}

De resulterende dataset houdt dan per cel bij: tijd, temperatuur, aantal bomen, aantal bomen in omgevende cellen, 
procentuele samenstelling van wegen en procentuele samenstelling van landgebruik in het algemeen.

\subsection{Analyse}

Met deze dataset kan analyse uitgevoerd worden op de relatie tussen temperatuur en boomplanting in verschillende stedelijke contexten in Gent.

De relatie tussen temperatuur en boomdensiteit wordt onderzocht met multiple lineaire regressie,
uitgevoerd in Python. Dit laat toe om het effect van boomdensiteit te kwantificeren terwijl gecontroleerd wordt 
voor andere variabelen zoals verhardingsgraad en landgebruik.


%---------- Verwachte resultaten ----------------------------------------------
\section{Verwacht resultaat, conclusie}%
\label{sec:verwachte_resultaten}

Op basis van de geplande methodologie wordt verwacht dat de multiple 
lineaire regressie een significant negatief verband zal aantonen tussen 
boomdensiteit en oppervlaktetemperatuur, in lijn met eerdere bevindingen 
uit \textcite{Knight2021}. De sterkte van dit effect zal naar verwachting 
variëren afhankelijk van lokale context: in sterk verharde gebieden wordt 
een groter koelend effect per boom verwacht dan in reeds groene wijken. 
Het regressiemodel zal deze effecten kwantificeren met p-waarden en 
prioriteitsgebieden identificeren waar vergroening het meest effectief 
zou zijn voor temperatuurreductie. Deze resultaten bieden een empirische 
onderbouwing voor lokaal klimaatadaptatiebeleid in Gent, 
waarbij beleidsmakers datagedreven beslissingen kunnen nemen 
over strategische vergroening om hittestress te mitigeren.
